\chapter{Konu Tanımı ve Senaryolar}

\section{Proje Amacı}
Bu proje, gerçek zamanlı sohbet uygulaması için mikroservis mimarisine dayalı bir yazılım sistemi geliştirmeyi amaçlamaktadır. Sistem; birebir ve grup sohbet özelliklerini destekleyen, dosya paylaşımı yapabilen, gerçek zamanlı mesajlaşma sağlayan, ölçeklenebilir ve modern bir mimari üzerine kurgulanmıştır.

\section{Proje Kapsamı}
Proje, aşağıdaki temel işlevleri kapsamaktadır:
\begin{itemize}
    \item Kullanıcı yönetimi ve kimlik doğrulama
    \item Birebir ve grup sohbet oluşturma
    \item Gerçek zamanlı mesajlaşma
    \item Mesaj geçmişi saklama ve görüntüleme
    \item Dosya ve medya paylaşımı
    \item Gerçek zamanlı bildirimler
\end{itemize}

Sistem en az 5 mikroservis içermektedir ve bu servisler bağımsız olarak geliştirilebilir, dağıtılabilir ve ölçeklendirilebilir şekilde tasarlanmıştır.

\section{Kullanıcı Senaryoları (User Stories)}

\subsection{Kullanıcı Yönetimi}
\begin{itemize}
    \item \textbf{US-1: Kullanıcı Kaydı:} Bir kullanıcı olarak, sisteme e-posta, kullanıcı adı ve şifre ile kayıt olabilmeliyim. Sistem e-posta benzersizliğini kontrol etmeli ve doğrulama e-postası göndermelidir.
    \item \textbf{US-2: Kullanıcı Girişi:} Kayıtlı hesabımla giriş yapabilmeli, sistem kimlik bilgilerimi doğrulamalı ve geçerli bir oturum token'ı döndürmelidir.
    \item \textbf{US-3: Profil Yönetimi:} Profil bilgilerimi (ad, soyad, resim, durum mesajı) görüntüleyebilmeli ve güncelleyebilmeliyim.
\end{itemize}

\subsection{Sohbet Yönetimi}
\begin{itemize}
    \item \textbf{US-4: Birebir Sohbet:} Başka bir kullanıcı ile birebir sohbet başlatabilmeliyim. Sistem mükerrer sohbet oluşturmamalıdır.
    \item \textbf{US-5: Grup Sohbeti:} Birden fazla kullanıcı ile grup sohbeti oluşturabilmeli, gruba isim verebilmeli ve katılımcı yönetebilmeliyim.
    \item \textbf{US-6: Sohbet Listesi:} Katıldığım tüm sohbetleri, son mesaj ve zaman bilgisine göre sıralı şekilde görebilmeliyim.
\end{itemize}

\subsection{Mesajlaşma}
\begin{itemize}
    \item \textbf{US-7: Mesaj Gönderme:} Sohbete metin mesajı gönderebilmeli ve mesajın iletildi/okundu durumlarını takip edebilmeliyim.
    \item \textbf{US-8: Mesaj Geçmişi:} Sohbetin mesaj geçmişini sayfalama (pagination) yapısı ile görüntüleyebilmeliyim.
    \item \textbf{US-9: Mesaj Silme:} Gönderdiğim mesajları silebilmeli, silinen mesajlar diğer kullanıcılara bildirilmelidir.
\end{itemize}

\subsection{Dosya Paylaşımı}
\begin{itemize}
    \item \textbf{US-10: Dosya Yükleme:} Sohbete resim, video veya doküman yükleyebilmeliyim. Sistem dosya türü ve boyut kontrolü yapmalıdır.
    \item \textbf{US-11: Dosya İndirme:} Paylaşılan dosyaları güvenli ve zaman sınırlı URL'ler üzerinden indirebilmeliyim.
\end{itemize}

\subsection{Bildirimler}
\begin{itemize}
    \item \textbf{US-12: Gerçek Zamanlı Bildirimler:} Yeni mesaj veya olaylarda anlık bildirim alabilmeliyim. Çevrimdışı isem bildirimler saklanmalı ve çevrimiçi olduğumda iletilmelidir.
\end{itemize}

\section{Ana Kullanım Akışları}

\subsection{Kullanıcı Kayıt ve Giriş Akışı}
\begin{enumerate}
    \item Kullanıcı kayıt sayfasına girer ve gerekli bilgileri doldurur.
    \item Sistem bilgileri doğrular, kullanıcıyı oluşturur ve doğrulama e-postası gönderir.
    \item Kullanıcı e-postasını doğrular ve giriş yapar.
    \item Sistem kimlik doğrulamasını yapar ve JWT token oluşturur.
    \item Kullanıcı ana sayfaya yönlendirilir.
\end{enumerate}

\subsection{Mesaj Gönderme Akışı}
\begin{enumerate}
    \item Kullanıcı sohbeti seçer ve mesajını yazar.
    \item İstemci mesajı \textit{Message Service}'e gönderir.
    \item Servis mesajı kaydeder ve \textit{Notification Service}'i tetikler.
    \item \textit{Notification Service}, \textit{WebSocket Gateway} üzerinden hedef kullanıcılara iletir.
    \item Çevrimiçi kullanıcılar mesajı anında görür, çevrimdışı kullanıcılar için kuyruğa alınır.
\end{enumerate}

\subsection{Dosya Paylaşımı Akışı}
\begin{enumerate}
    \item Kullanıcı dosya yükleme isteği gönderir.
    \item \textit{File Service} dosyayı doğrular ve depolama alanına kaydeder.
    \item Dosya için erişim URL'i oluşturulur.
    \item İstemci, bu URL'i içeren bir mesaj oluşturarak sohbete gönderir.
\end{enumerate}
